% =====================================================================
% Working note
% Is the signal channel free-space, and how to add the body to h
% Companion to the coherent-exposure-operator paper (paper C),
% q_complement.tex, and unified_physics.tex
% =====================================================================

\documentclass[11pt]{article}

\usepackage[a4paper,margin=2.6cm]{geometry}
\usepackage[utf8]{inputenc}
\usepackage[T1]{fontenc}
\usepackage{lmodern}
\usepackage{microtype}
\usepackage{amsmath,amssymb,amsthm,mathtools}
\usepackage{bm}
\usepackage{booktabs}
\usepackage{xcolor}
\usepackage[colorlinks=true,allcolors=blue!50!black]{hyperref}
\usepackage[capitalize]{cleveref}

% Macros matching paper C (build/main.tex) and q_complement.tex
\newcommand{\khat}{\hat{\bm{k}}}
\newcommand{\nhat}{\hat{\bm{n}}}
\newcommand{\rr}{\mathbf{r}}
\newcommand{\xx}{\mathbf{x}}
\newcommand{\hh}{\mathbf{h}}
\newcommand{\GG}{\mathbf{G}}
\newcommand{\Gtil}{\tilde{\mathbf{G}}}
\newcommand{\Gre}{\Gtil^{\mathrm{re}}}
\newcommand{\Qmat}{\mathbf{Q}}
\newcommand{\Qabs}{\mathbf{Q}_{\mathrm{ab}}}
\newcommand{\Qref}{\mathbf{Q}_{\mathrm{re}}}
\newcommand{\Jmat}{\mathbf{J}}
\newcommand{\Phimat}{\mathbf{\Phi}}
\newcommand{\Psimat}{\bm{\Psi}}
\newcommand{\Fn}{\mathbf{F}_n}
\newcommand{\Rn}{\mathbf{R}_n}
\newcommand{\psin}{\bm{\psi}_n}
\newcommand{\diff}{\mathrm{d}}
\newcommand{\Complex}{\mathbb{C}}
\newcommand{\ntilde}{\tilde{n}}
\newcommand{\Vis}{V}
\newcommand{\rue}{\rr_{\mathrm{UE}}}
\newcommand{\hb}{\hh_{\mathrm{b}}}
\newcommand{\uue}{\mathbf{u}_{\mathrm{UE}}}
\newcommand{\mfock}{m_{\mathrm{F}}}
\DeclareMathOperator{\RE}{Re}
\DeclareMathOperator{\diag}{diag}

\theoremstyle{remark}
\newtheorem{remark}{Remark}

\title{Is the signal channel free-space, and how to add the body to $\hh$}
\author{Working note, companion to paper C and \texttt{q\_complement}}
\date{2026-08-10}

\begin{document}
\maketitle

\begin{abstract}
\noindent
The signal channel $\hh$ of paper C is the free-space channel of a
scene that does not contain the body. The body enters the model only
through the exposure branch. This note makes that statement precise,
argues why it matters for the MRT baseline and for the small
device-to-body separations of the sweep, and then assembles a
body-perturbed channel $\hb$ without introducing any new physics
object. The scattering correction is the emitted-field channel of
\texttt{q\_complement}, the same $\Rn$-built channel whose Gram is
$\Qref$, sampled at the device instead of integrated over exit
directions: $\Qref$ is what the body emits in total, $\Delta\hh$ is
what the device hears of that emission. Lorentz reciprocity turns the
sample into a reaction integral on the body surface between the
reflected channel and the UE antenna's spherical-wave field, so no
specular-point search and no far-field assumption on the body are
needed. The shadowing correction is the visibility gate of the
exposure branch evaluated at the device, with the Fock transition
function of the unified-physics note as its matched refinement. To
first order in the body interaction $\Qmat$ is unchanged, so the
consistent upgrade is: replace $\hh$ by $\hb$ in the precoder design,
keep $\Qmat$.
\end{abstract}

\section{Where the body enters the model, and where it does not}

Paper C builds three linear maps from one set of ray-traced path
triples $(j(n), \khat_n, \psin)$:
\begin{center}
\begin{tabular}{@{}llll@{}}
\toprule
Branch & Map & Amplitude block & Body present \\
\midrule
Free-space field & $\mathbf{E}(\rr) = \GG(\rr)\,\xx$,
  \ $\GG = \Psimat\,\Phimat(\rr)\,\Jmat$ & $\Psimat$ & no \\
Signal & $y = \hh^T \xx$,
  \ $\hh^T = \bm{\alpha}^T\,\Phimat(\rue)\,\Jmat$ & $\bm{\alpha}^T$ & no \\
Exposure & $S_{\mathrm{ab}}(\rr) = \|\Gtil(\rr)\,\xx\|^2$,
  \ $\Gtil = \tilde{\bm{\Psi}}(\rr)\,\Phimat(\rr)\,\Jmat$ & $\tilde{\bm{\Psi}}(\rr)$ & yes \\
\bottomrule
\end{tabular}
\end{center}
The ray trace runs in a scene without the phantom. The body is
inserted a posteriori, on the exposure branch only, as a surface that
intercepts the free-space paths and transmits them through the Fresnel
operator $\Fn(\rr)$. The signal branch samples the same free-space
field at the device position $\rue$ through the UE antenna response,
so $\hh$ is the channel of a world in which the body does not exist.
Two consequences follow.

First, the body does not shadow the device. A path arriving at $\rue$
from behind the phantom, through the torso, is counted at full
strength in $\hh$. Measurement campaigns at millimeter-wave
frequencies put human-body blockage at roughly 20 to 40\,dB, so these
are not small terms.

Second, the body does not scatter into the device. The sweep of
paper C moves the served point from 0 to 20\,cm off the chest, and
states that the device separation enters only through $\hh$. In the
model that is true because $\hh$ ignores the chest entirely. In
reality a device a few centimeters from a high-index dielectric
half-space sits in a two-ray standing wave, and its received amplitude
ripples with separation. Both effects perturb the direction of
$\hh^{*}$, which is the MRT precoder, which is the baseline against
which ECBF is scored, and which sets the exposure
$\xx^H \Qmat\, \xx$ itself.

\begin{remark}[The measured channel contains the body]
In a deployed system $\hh$ is not computed, it is estimated from
pilots that physically propagate past the user's body. Real CSI is
body-inclusive by construction, and real MRT is matched to the
body-inclusive channel. The free-space $\hh$ is therefore not the
conservative choice, it is a channel no receiver would ever measure.
Adding the body to $\hh$ moves the simulation toward the channel the
network actually sees.
\end{remark}

\section{Mechanisms and expected size}

\paragraph{Shadowing.}
In the LOS geometry of paper C the phantom stands behind the served
device, so environment paths that bounce off the far half of the hall
arrive at $\rue$ through the torso. Free-space MRT co-phases those
paths and spends transmit power on them. That power is wasted for the
link, and it is precisely the power that terminates in the trunk,
since the body is what blocks it. A body-aware $\hh$ suppresses these
paths, and MRT on it stops illuminating the body from behind. Some of
the exposure reduction that paper C attributes to ECBF may already be
delivered by matched filtering on the correct channel. That is a
hypothesis worth one experiment, not a claim.

\paragraph{Chest reflection.}
At normal incidence on skin the field reflection coefficient is
$\Gamma = (1 - \ntilde)/(1 + \ntilde)$, and with
$|\ntilde| \approx 4.83$ at 28\,GHz, $|\Gamma| \approx 0.66$. A device
at separation $d$ on the chest normal receives direct and
chest-reflected copies of the dominant path, with relative phase
$2 k_0 d$, so the received amplitude of that path is modulated by
$|1 + D\,\Gamma\, e^{-2 i k_0 d}|$ with $D \le 1$ the divergence
factor of the curved chest. On a flat chest the swing is $+4.4$ to
$-9.3$\,dB. Curvature reduces it, but at $d \le 3$\,cm the specular
patch is small and $D$ is close to one, so several dB of ripple in
$|\hh|$ across the separation sweep is the expectation, not the
exception. The current model shows none of it.

\section{The body as an emitter, and the device in its field}
\label{sec:emitter}

The elegant route does not build a new scattering model. It reuses
the emitted-field channel that \texttt{q\_complement} already defines
for power bookkeeping, and evaluates it at one more point.

\subsection{The outgoing channel, recalled}

\texttt{q\_complement} introduces the Fresnel reflection operator,
sibling of $\Fn$, with coefficients
\begin{equation}\label{eq:r-coef}
  r_{s,n} = \frac{\mu_n - \xi_n}{\mu_n + \xi_n},
  \qquad
  r_{p,n} = \frac{\ntilde^2 \mu_n - \xi_n}{\ntilde^2 \mu_n + \xi_n},
\end{equation}
linked to the transmission coefficients by $t_{s} = 1 + r_{s}$ and
$\ntilde\, t_{p} = 1 + r_{p}$. At the field level, for a front-facing
path at $\rr \in \Sigma$, the operator that maps the incident
polarization to the outgoing surface field is
\begin{equation}\label{eq:Rn}
  \Rn(\rr) =
  \bigl(r_{s,n}\, \hat{e}_{s,n} \hat{e}_{s,n}^{T}
  + r_{p,n}\, \hat{e}_{p,n}^{\mathrm{r}} \hat{e}_{p,n}^{T}\bigr)\,
  \Vis(\rr, \khat_n)\, H(\mu_n),
  \qquad
  \hat{e}_{p,n}^{\mathrm{r}} = \hat{e}_{s,n} \times \khat_n^{\mathrm{r}},
\end{equation}
with $\khat_n^{\mathrm{r}} = \khat_n + 2 \mu_n \nhat$ the specular
direction. One refinement relative to \texttt{q\_complement}: its
$\Rn$ keeps the incident TM basis because only $|\Rn \psin|^2$ enters
the power operator $\Qref$. For a field evaluated off the surface the
outgoing polarization direction matters, hence the rotated basis
$\hat{e}_{p,n}^{\mathrm{r}}$, and the visibility gate that
\texttt{q\_complement} leaves implicit is written out.

The reflected surface field under precoder $\xx$ is then the outgoing
channel
\begin{equation}\label{eq:Gre}
  \mathbf{E}^{\mathrm{re}}(\rr) = \Gre(\rr)\,\xx,
  \qquad
  \Gre(\rr) = \bigl[\Rn(\rr)\psin\, e^{-i k_0 \khat_n \cdot \rr}\bigr]_{n}
  \, \Jmat \in \Complex^{3 \times M},
\end{equation}
the exact parallel of $\Gtil$ with $\Rn$ in place of $\Fn$. This is
the field the body emits back into the environment. It is one object
with three readings, and the third is the new one:
\begin{center}
\begin{tabular}{@{}lll@{}}
\toprule
Reading & Operation on $\Gre$ & Object \\
\midrule
Total re-emission & Gram, integrated over $\Sigma$ & $\Qref$ \\
Thermal dual & reciprocity, array as receiver & Kirchhoff identity of \texttt{q\_complement} \\
Device coupling & Huygens sample at $\rue$ & $\Delta\hh$, this note \\
\bottomrule
\end{tabular}
\end{center}
$\Qref$ answers ``how much does the body emit in total''. The
correction to the signal channel answers ``how much of that emission
does the device hear''. Same surface response, read at a point
instead of in aggregate.

\subsection{The device sample, in reaction form}

The Huygens continuation of the aperture field \cref{eq:Gre} to an
exterior point gives the body-scattered field at the device,
schematically
$\mathbf{E}_{\mathrm{sc}}(\rue) = \int_\Sigma \mathbf{K}(\rue, \rr)\,
\Gre(\rr)\,\xx\, \diff A$ with $\mathbf{K}$ the free-space propagation
kernel. Projecting on the UE antenna and applying Lorentz reciprocity
folds the kernel and the antenna response into a single physical
object, the field the UE antenna itself radiates on the body surface.
Let $\uue(\rr) \in \Complex^3$ be that field under unit port
excitation, the exact spherical wave of the device antenna, all
$1/(k_0 s)$ orders retained. The coupling is the reaction integral
\begin{equation}\label{eq:reaction}
  \Delta\hh^{T} = \mathbf{a}^{T} \Jmat,
  \qquad
  a_n = \kappa \int_{\Sigma}
  \uue(\rr)^{T}\, \Rn(\rr)\, \psin\,
  e^{-i k_0 \khat_n \cdot \rr}\, \diff A,
\end{equation}
with $\kappa$ a reciprocity constant fixed by the port normalization
conventions, to be pinned at implementation time. The structure is
the point: $\Delta\hh$ is a bilinear overlap on $\Sigma$ between the
array-side outgoing channel and the device-side illumination. The
body-perturbed signal channel is
\begin{equation}\label{eq:hb}
  \boxed{\;
  \hb^{T}
  = \bm{\alpha}^{T}\, \mathbf{V}_{\mathrm{UE}}\, \Phimat(\rue)\, \Jmat
  \;+\; \mathbf{a}^{T} \Jmat ,
  \;}
\end{equation}
with $\mathbf{V}_{\mathrm{UE}}$ the device-side shadowing gate of
\cref{sec:shadow}. Three properties make \cref{eq:reaction} the right
form.

\emph{No stationary-phase machinery.} The integral is evaluated by
the same triangle quadrature that assembles $\Qmat$, on the same
per-triangle quantities: incidence bases, Fresnel coefficients,
visibility, plus one new spherical-wave evaluation per triangle per
device position. No specular-point search on the mesh, no curvature
radii, no divergence factor. Geometrical optics reappears only as
the asymptotic evaluation: the phase of \cref{eq:reaction} is
$-k_0(\khat_n \cdot \rr + s)$ with $s = |\rue - \rr|$, its stationary
point is the specular point of path $n$, and the stationary-phase
value is the GO reflected ray with the divergence factor of the local
curvature. That is the sanity check, not the method.

\emph{No far-field assumption on the body.} The question of whether
the device is in the far field of the body never arises, because the
body is never treated as a single radiator. Each surface element
couples to the device through the exact $\uue$, which is valid from
the radiating near field outward. The regime boundaries of the
unified-physics note apply on the device side only: the reactive
boundary $\lambda/2\pi \approx 1.7$\,mm at 28\,GHz sits far below the
1\,cm separation grid, so every nonzero separation of the sweep is in
the spherical-wave regime, and the spherical wave costs nothing. The
single excluded case is contact, $d = 0$, where reactive coupling and
antenna detuning dominate and no surface method applies.

\emph{Linear in $\xx$, factorization-compatible.} Equation
\cref{eq:reaction} keeps the amplitude-times-dispatch shape of every
channel in paper C. The phase is referenced on $\Sigma$ rather than
at $\rue$, so the term does not factor through $\Phimat(\rue)$, which
is harmless since the precoder design consumes $\hb$ as one
$M$-vector.

\paragraph{Two-ray sanity check.}
For the dominant path at normal incidence on the chest and the device
on the chest normal at separation $d$, the stationary point of
\cref{eq:reaction} sits directly behind the device, and the
stationary-phase evaluation gives
$a_n = D\, \Gamma\, e^{-2 i k_0 d}\, \alpha_n\,
e^{-i k_0 \khat_n \cdot \rue}$ with $\Gamma$ from \cref{eq:r-coef} at
$\mu_n = 1$ and $D$ the chest divergence factor. The separation sweep
of paper C acquires exactly the textbook two-ray ripple.

\section{Device-side shadowing, with the Fock gate}
\label{sec:shadow}

The exposure channel gates each path by the binary visibility
$\Vis(\rr, \khat_n)$. The same function, evaluated at the device with
the phantom as occluder, gives
\begin{equation}\label{eq:Vue}
  \mathbf{V}_{\mathrm{UE}}
  = \diag\bigl(g_1, \ldots, g_N\bigr),
  \qquad
  g_n = \Vis(\rue, \khat_n) \in \{0, 1\},
\end{equation}
one BVH ray query per path, the identical primitive the
ambient-occlusion bake uses. The binary gate is the level-consistent
default, matching the fidelity at which the exposure branch runs.

The refinement is not knife-edge diffraction. The occluders are
limbs and the torso, convex on the wavelength scale, and the
unified-physics note shows that the correct shadow-boundary physics
on a convex body is Fock theory: the binary gate is replaced by the
complex transition function $g(\xi_n)$ of the Fock boundary layer,
with detour parameter $\xi_n = \mfock\,\theta_n$,
$\mfock = (k_0 R_{\mathrm{F}}/2)^{1/3}$, $R_{\mathrm{F}}$ the
in-plane curvature radius of the occluding contour, and $\theta_n$
the signed grazing angle of the blocked ray past that contour. The
soft and hard polarization split and the lossy-skin impedance
correction carry over unchanged, since the gate is the same object
the exposure branch uses at its own shadow edges. The penumbra width
scales as $(k_0 R_{\mathrm{F}})^{-1/3}$, a few degrees for an arm at
28\,GHz, which is exactly the angular scale the separation and
position sweeps of paper C step through, so the smooth gate matters
for the gradient of $\hb$ along the sweep even where the binary gate
gets the mean right. Fidelity pairing is the rule: binary with
binary, Fock with Fock, on both branches.

\section{What happens to the exposure operator}

Nothing, to first order, and that is the point. Order the model by
powers of the surface interaction. The exposure branch is first
order, one $\Fn$ per path. The signal branch of paper C is zeroth
order, the body absent. The two corrections of \cref{eq:hb} promote
the signal branch to first order, one $\Rn$ or one occlusion event
per path. A body correction to $\Qmat$ itself requires two surface
interactions, reflect off one body part and absorb on another, which
is the inter-body recapture of the unified-physics note and the first
term of the Neumann cascade of \texttt{q\_complement}. On a mostly
convex body it is geometrically starved, active in the concavities,
and the unified-physics note already keeps it off by default because
the body-averaged enhancement sits inside the diffraction and
tissue-property error bars. The consistent first-order theory is
therefore
\begin{equation}
  \xx^{\star}
  \propto (\lambda \Qmat + \nu \mathbf{I})^{-1} \hb^{*},
  \qquad \Qmat \text{ unchanged},
\end{equation}
and likewise MRT becomes $\xx \propto \hb^{*}$. Every downstream
quantity of paper C, the Pareto sweep, the binding-restriction
analysis, the concentration statistics, is a function of
$(\hh, \Qmat)$ and runs unmodified on $(\hb, \Qmat)$.

\begin{remark}[One quadrature, three operators]
After this note the model evaluates three surface functionals on the
same mesh quadrature: $\int \Gtil^H \Gtil$ for $\Qabs$,
$\int (\Gre)^H \Gre$ for $\Qref$ when the bookkeeping is wanted, and
$\int \uue^T \Rn \psin\, e^{-i k_0 \khat_n \cdot \rr}$ for
$\Delta\hh$. The per-triangle geometry, Fresnel bases, and
visibility are shared. The marginal cost of the body-aware channel
is one spherical-wave evaluation per triangle per device position,
$O(N_\triangle)$, negligible against the exposure kernel.
\end{remark}

\section{Implementation routes and first measurements}

\paragraph{Closed form, from existing machinery.}
Shadowing: $N$ BVH queries per device position. Scattering: the
reaction quadrature above. The only genuinely new code is $\uue$,
the exact field of the device dipole on the mesh, and the reciprocity
constant $\kappa$, which should be pinned by one validation case, for
example a flat phantom slab against the analytical two-ray answer.

\paragraph{Body-in-scene trace, as the oracle.}
Sionna RT accepts the phantom mesh with a tissue material, and a
trace of that scene returns a body-inclusive channel directly,
including the multi-bounce body-environment interactions that
\cref{eq:hb} truncates. Two cautions. First, double counting: the
exposure operator must keep being fed by the body-free trace, since
$\Fn$ models the surface interaction analytically. The clean
protocol is two traces of the same scene, body-in for $\hh$, body-out
for $\Qmat$. Second, cost: the body-in trace repeats per posture and
position, which is the per-configuration expense the closed-form
operator exists to avoid. The oracle validates \cref{eq:reaction} on
a few configurations, it does not replace it in the sweep.

\paragraph{What to measure first.}
Three numbers decide whether this matters. One, the fraction of
$\|\hh\|^2$ carried by torso-blocked paths across the sweep, from
$\mathbf{V}_{\mathrm{UE}}$ alone with no scattering machinery. Two,
the ripple of $\|\hb\|$ over the 1 to 3\,cm separations, from the
reaction quadrature, checked against the two-ray formula. Three, the
change in $\xx^H \Qmat\, \xx$ at fixed received signal power when MRT
switches from $\hh^{*}$ to $\hb^{*}$, whose sign is genuinely open,
since shadowing steers power away from the trunk while the chest
reflection rewards paths that bounce off the chest into the device.

\end{document}
